\section{The Kick}
\label{sec:kick}

We compute the kick that ring $k$ receives from the rings inside it,
for $1\le k<n$. While ring $k$ is quiet, the inner field $\hin_k$ of
Definition~\ref{def:regions} drives the unstable coefficient of ring $k$ exactly
through the mode-$3$ moment $q_{k+1}$ of the region $D_{k+1}(t)$: the unstable
coefficient of the part of the linear response of ring $k$ driven by $\hin_k$
solves the scalar linear Equation~\eqref{eq:kick}, with drive $q_{k+1}$
(Lemma~\ref{lem:coupling}). Proposition~\ref{prop:kick} bounds the modulus of
the solution at the final time, with its free growth factored out, from above
and below, and Section~\ref{sec:relay} uses these bounds to start ring $k$
(Lemmas~\ref{lem:handoff} and~\ref{lem:kicktrue}).

\begin{lemma}
\label{lem:coupling}
Assume \textup{(N1)}--\textup{(N3)}, fix $n\ge2$ and the objects of
Definitions~\ref{def:regions} and \ref{def:clock}, and let $1\le k<n$.
\begin{enumerate}[label=\textup{(\arabic*)}]
\item Let $t\ge0$ with $D_{k+1}(t)\subset B_{ar_k/8}$, and put $\tau=A_kt$. Then
  $-\hin_k(\tau)\cdot\nabla\Phi$, extended by $0$ outside $N'$, lies in $X^s$
  for every $s\ge0$, and
  \[
  \ellst(G_k)\big(-\hin_k(\tau)\cdot\nabla\Phi\big)=\mu\,\beta_k\,q_{k+1}(t).
  \]
  Moreover, with $\nu^{\mathrm{in}}_k$ of Definition~\ref{def:regions} at the
  time $t$, the mode-$3$ profile of $\hin_k(\tau)\cdot\er$ is
  $\frac i2\pout_3[\nu^{\mathrm{in}}_k]\,r^{-4}$ for $a/4\le r\le3$.
\item Assume that the clock is defined up to level $k+1$ and that
  $D_{k+1}(t)\subset B_{ar_k/8}$ for every $t\in[0,\tsh_k]$. Then
  $\tau\mapsto-\hin_k(\tau)\cdot\nabla\Phi$ is continuous from $[0,\tash_k]$ into
  $X^s$ for every $s$, and we may put, in ring-$k$ time,
  \[
  \wlin_{\mathrm{in},k}(\tau)=\int_0^\tau e^{(\tau-v)\Lin(G_k)}\big(-\hin_k(v)\cdot\nabla\Phi\big)\dd v
  \qquad(0\le\tau\le\tash_k).
  \]
  Regard $q_{k+1}$ and $\cin=\ellst(G_k)\big(\wlin_{\mathrm{in},k}\big)$ as functions
  of the kick time $\tauK=A_{k+1}(t-T_{k+1})\in[\tau_0,\Tpin]$. Then
  $\cin\in C^1([\tau_0,\Tpin])$ and, exactly,
  \begin{equation}
  \label{eq:kick}
  \frac{d\cin}{d\tauK}=\zeta_0\,\cin+\beta_k\,q_{k+1},\qquad \zeta_0=\frac{\lambda_k}{\mu},\qquad
  \cin(\tau_0)=0 .
  \end{equation}
  The \textbf{kick amplitude} $\Ain(\tauK)=e^{-\zeta_0\tauK}\cin(\tauK)$ satisfies
  \[
  \cin_k:=\ellst(G_k)\big(\wlin_{\mathrm{in},k}(\tash_k)\big)=e^{\zeta_0\Tpin}\Ain(\Tpin).
  \]
\item If, in addition to the hypotheses of \textup{(2)}, the rings $i<k$ lie
  outside $B_{6r_k}$ for every $t\in[0,\tsh_k]$, then
  $\tau\mapsto-\hout_k(\tau)\cdot\nabla\Phi$ is continuous from $[0,\tash_k]$ into
  $X^s$ for every $s$, and $\wlin_{\mathrm{out},k}$, defined as
  $\wlin_{\mathrm{in},k}$ with $\hout_k$ in place of $\hin_k$, satisfies $\wlin_{\mathrm{in},k}+\wlin_{\mathrm{out},k}=\wlin$ on
  $[0,\tash_k]$, where $\wlin$ is the linear response of
  Theorem~\ref{thm:quiet} for ring $k$, with $G=G_k$ and $h=h_k$.
\end{enumerate}
\end{lemma}
\begin{proof}
\emph{(1).} This part makes precise the computation made before
Definition~I.\ref{I-def:eigendata}, where $b$ is introduced as the mode-$3$
forcing of the ring by an interior field of unit outer moment, and $\beta(G)$ as
the component of that forcing along the unstable eigenfunction. By
Lemma~\ref{lem:regions-basic}(2),(3), the function
$\nu^{\mathrm{in}}_k$ of Definition~\ref{def:regions} at time $t$, the vorticity of
$D_{k+1}(t)$ in ring-$k$ units, is real,
$3$-fold rotationally symmetric, in $C_c^\infty(\R^2)$, and its support avoids $0$; by the
hypothesis $D_{k+1}(t)\subset B_{ar_k/8}$ the support is a compact subset of
$B_{a/8}$, so the largest modulus $\rho$ of a point of the support satisfies $\rho<a/8$. By
Lemma~I.\ref{I-lem:BS-basic}(d), $\KBS\ast\nu^{\mathrm{in}}_k=\nabla^\perp\psi$ with
$\psi=\frac1{2\pi}\log|\cdot|\ast\nu^{\mathrm{in}}_k$ smooth on $\{|y|>\rho\}$.
The monopole $G_ky^\perp/|y|^2$ is orthogonal to $\er$, so the polar form
$v\cdot\er=-r^{-1}\partial_\theta\psi$ of the Biot-Savart law
(Section~I.\ref{I-sec:prelim}) gives $\hin_k(\tau)\cdot\er=-r^{-1}\partial_\theta\psi$
on $N'$. Put $p_3=\pout_3[\nu^{\mathrm{in}}_k]$, the outer mode-$3$ moment of
Definition~I.\ref{I-def:multipole}. By the expansion of the stream function stated
before Lemma~\ref{lem:ext-multipole}, with $m=3$ and $\varrho=\rho$, the angular mode $3$ of $\psi$ is
$-\frac16p_3r^{-3}$ for $r>\rho$, and by integration by parts in $\theta$ that of
$\partial_\theta\psi$ is $3i$ times that of $\psi$. Since
$N'=\{a/4\le|y|\le3\}\subset\{|y|>\rho\}$, the mode-$3$ profile of
$\hin_k(\tau)\cdot\er$ is $\frac i2p_3r^{-4}$ for $a/4\le r\le3$, the last claim
of (1). As $\Phi$ is radial and $\Phi'$ vanishes outside $[a,1]$,
$\hin_k(\tau)\cdot\nabla\Phi=\Phi'(r)\,\hin_k(\tau)\cdot\er$, so the mode-$3$
profile of $-\hin_k(\tau)\cdot\nabla\Phi$ is
$-\frac i2p_3r^{-4}\Phi'(r)=p_3\,b(r)$, with $b$ of Definition~I.\ref{I-def:eigendata}.
By Equation~\eqref{eq:BS-complex}, $\hin_k(\tau)$ is the restriction to
$N'$ of a field smooth on $\{|y|>\rho\}$, and it is $3$-fold rotationally symmetric because
$\nu^{\mathrm{in}}_k$ is; so $-\hin_k(\tau)\cdot\nabla\Phi$ is real, smooth,
$3$-fold rotationally symmetric and supported in $\{a\le|y|\le1\}\subset N$, hence in $X^s$
for every $s$. By Definition~I.\ref{I-def:eigendata}, $\ellst(G_k)(w)=n_*(G_k)\,\ellst(w_3)$
depends linearly on the mode-$3$ profile $w_3$ of $w$ only, and
$\beta(G_k)=\ellst(G_k)(b\,e^{3i\theta})=n_*(G_k)\ellst(b)$; so
$\ellst(G_k)(-\hin_k(\tau)\cdot\nabla\Phi)=p_3\,\beta(G_k)$. It remains to
express $p_3$ through the normalized moment $q_{k+1}(t)$ of
Definition~\ref{def:regions}. The change
of variables $x=r_ky$ and the reality of $\omega_{\ge k+1}(t)$ give
\[
p_3=\frac1{2\pi A_k}\int\omega_{\ge k+1}(t,r_ky)\,\bar y^{\,3}\dd y
=\frac{\overline{\Qm[\omega_{\ge k+1}(t)]}}{2\pi A_kr_k^5}
=\frac{A_{k+1}r_{k+1}^5}{2\pi A_kr_k^5}\,q_{k+1}(t)=\frac{\mu s_k^5}{2\pi}\,q_{k+1}(t),
\]
and (1) follows from $m+2=5$ (recall $m=3$) and $\beta_k=\beta(G_k)s_k^5/(2\pi)$
(Definition~\ref{def:regions}).

\emph{(2).} By the hypothesis and the paragraph after
Lemma~\ref{lem:regions-basic}, $\tau\mapsto\hin_k(\tau)$ is continuous from
$[0,\tash_k]$ into $C^M(N')$ for every integer $M\ge0$. Put
$g(\tau)=-\hin_k(\tau)\cdot\nabla\Phi$. By (1), $g(\tau)\in X^s$ for every $s$,
and by Lemma~I.\ref{I-lem:uniform},
$\|g(\tau)-g(\tau')\|_{H^M}\le\Cnab(M)\|\hin_k(\tau)-\hin_k(\tau')\|_{C^M(N')}$;
so $g$ is continuous into $X^s$ for every $s$. Since $e^{t\Lin(G_k)}$ is a
$C_0$-semigroup on $X^s$ (Lemma~I.\ref{I-lem:semigroup-construct}), the integrand
$v\mapsto e^{(\tau-v)\Lin(G_k)}g(v)$ is continuous into $X^s$, so
$\wlin_{\mathrm{in},k}$ is well defined.

Next we change from ring-$k$ time to kick time. Recall from
Definition~\ref{def:clock} that $\tau_0=-A_{k+1}T_{k+1}$ is the kick time of $t=0$.
Since $A_{k+1}=\mu A_k$ by Equation~\eqref{eq:datum}, we have
$A_kT_{k+1}=-\tau_0/\mu$, so ring-$k$
time and kick time are related by $\tau=(\tauK-\tau_0)/\mu$. As
$\tsh_k=T_{k+1}+\Tpin/A_{k+1}$ (Definition~\ref{def:clock}), $\tauK$ runs over
$[\tau_0,\Tpin]$ as $t$ runs over $[0,\tsh_k]$. As $\ellst(G_k)$ is bounded on
$X^0$ (Lemma~I.\ref{I-lem:uniform}), it commutes with the integral that defines
$\wlin_{\mathrm{in},k}$. In the display below, the first equality is then the
last identity of Proposition~I.\ref{I-prop:semigroup}, by which $\ellst(G_k)$
turns $e^{t\Lin(G_k)}$ into multiplication by $e^{\lambda_kt}$; the second is
(1); and the third is the substitution
$v=(\tauK'-\tau_0)/\mu$. This gives, for $\tau_0\le\tauK\le\Tpin$ and
$\tau=(\tauK-\tau_0)/\mu$, with $q_{k+1}$ written as a function of kick time in
the last line,
\begin{align*}
\cin(\tauK)&=\int_0^\tau e^{\lambda_k(\tau-v)}\,\ellst(G_k)\big(g(v)\big)\dd v
=\mu\beta_k\int_0^\tau e^{\lambda_k(\tau-v)}\,q_{k+1}(v/A_k)\dd v\\
&=\beta_k\int_{\tau_0}^{\tauK}e^{\zeta_0(\tauK-\tauK')}\,q_{k+1}(\tauK')\dd\tauK'.
\end{align*}
By Equation~\eqref{eq:omega-ge} and area preservation,
$\Qm[\omega_{\ge k+1}(t)]=\int\sum_{i>k}f^{(n)}_i(x')\,\flow_t(x')^m\dd x'$. The
integrand is supported in the compact set $D_{k+1}(0)$, it is bounded uniformly
for $t$ in compact intervals, and it is continuous in $t$, as
$(t,x')\mapsto\flow_t(x')$ is continuous (Theorem~\ref{thm:smooth}\textup{(4)});
so, by dominated convergence, $q_{k+1}$ is continuous. Hence, taking the factor
$e^{\zeta_0\tauK}$ out of the last integral and using the fundamental theorem of
calculus, $\cin\in C^1([\tau_0,\Tpin])$, $\cin(\tau_0)=0$, and differentiation gives
Equation~\eqref{eq:kick}. At $t=\tsh_k$ we have $\tau=\tash_k$ and
$\tauK=A_{k+1}(\tsh_k-T_{k+1})=\Tpin$, so $\cin_k=\cin(\Tpin)=e^{\zeta_0\Tpin}\Ain(\Tpin)$.

\emph{(3).} By the hypotheses and the paragraph after
Lemma~\ref{lem:regions-basic}, $\tau\mapsto\hout_k(\tau)$ is continuous from
$[0,\tash_k]$ into $C^M(N')$ for every $M$. By Lemma~\ref{lem:regions-basic}(3),
$\nu^{\mathrm{out}}_k$ is real and $3$-fold rotationally symmetric, and, as the
rings $i<k$ lie outside $B_{6r_k}$, it vanishes on $B_6$; so
$\hout_k(\tau)$ is real, $3$-fold rotationally symmetric and smooth on $B_6$, by
Equation~\eqref{eq:BS-complex}, and $B_6$ contains $N'$. Hence, as in (1),
$-\hout_k(\tau)\cdot\nabla\Phi\in X^s$
for every $s$, and it is continuous in $\tau$ into $X^s$ by
Lemma~I.\ref{I-lem:uniform}, as in (2). By the discussion after
Lemma~\ref{lem:regions-basic}, ring $k$ in ring-$k$ units satisfies the
hypotheses of Lemma~\ref{lem:reduction} on $[0,\tash_k]$ with $G=G_k$ and
$h=h_k=\hin_k+\hout_k$, so the linear response $\wlin$ of
Theorem~\ref{thm:quiet}\textup{(b)} is defined with this $h$; the identity for
$\wlin$ is then the linearity of the integral that defines it.
\end{proof}

Proposition~\ref{prop:kick} works in the following setting, which isolates
Equation~\eqref{eq:kick}. Let $\zeta_0,\beta\in\C$, put $x=\Re\zeta_0$, let
$\tau_0<\tau_1$ be real, put $I=[\tau_0,\tau_1]$, let
$q\colon I\to\C$ be continuous, let $c$ solve $c'=\zeta_0c+\beta q$ on $I$ with
$c(\tau_0)=0$, and put $A(\tau)=e^{-\zeta_0\tau}c(\tau)$. The factor
$e^{-\zeta_0\tau}$ removes the free growth of $c$, so that $A$ is a weighted
integral of the drive $q$, as the next display shows. Since
$A'=e^{-\zeta_0\tau}(c'-\zeta_0c)=\beta e^{-\zeta_0\tau}q$ and $A(\tau_0)=0$,
\begin{equation}
\label{eq:kick-A}
A(\tau)=\beta\int_{\tau_0}^\tau e^{-\zeta_0s}\,q(s)\dd s\qquad(\tau\in I).
\end{equation}
Under the hypotheses of Lemma~\ref{lem:coupling}\textup{(2)}, the functions
$\cin$, $q_{k+1}$ and $\Ain$, with $\zeta_0=\lambda_k/\mu$, $\beta=\beta_k$ and
$\tau_1=\Tpin$, are an instance of this setting; here
$\tau_0=-A_{k+1}T_{k+1}<0<\Tpin$ by Lemma~\ref{lem:clock}(1) and
Definition~\ref{def:clock}.

Recall $z=\lamst(0)$ and $\sigma=\Re z>0$ (Definition~I.\ref{I-def:eigendata},
Proposition~I.\ref{I-prop:U}(a)), and the integral $\Ip{p}{x}$ of
Equation~(I.\ref{I-eq:Ip}), a decreasing function of $x$.

\begin{proposition}
\label{prop:kick}
Assume \textup{(N1)}--\textup{(N3)}, work in the setting above, and let
$1<\mu\le2$ and $G,G'\in[0,\Gone]$ with $\zeta_0=\lamst(G)/\mu$ and
$|\lamst(G)-z|,|\lamst(G')-z|\le\sigma(\mu-1)/(8\mu)$. Put $z'=\lamst(G')$ and
$w=z'-\zeta_0$. Let $T_g>0$ with $\tau_0<-T_g$, assume $\tau_1>0$, let $\bar Q>0$
and let $p\ge0$ be an integer. Assume:
\begin{description}[font=\normalfont]
\item[\textup{(Q1)}] there are $\qs\in\C\setminus\{0\}$ and a measurable
  $\err\colon[-T_g,0]\to[0,\infty)$ with $|q(\tau)-\qs e^{z'\tau}|\le\err(\tau)$ on
  $[-T_g,0]$;
\item[\textup{(Q2)}] there is a measurable function
  $Q_{\mathrm{pre}}\colon[\tau_0,-T_g]\to[0,\infty)$ with
  $|q(\tau)|\le Q_{\mathrm{pre}}(\tau)$ on $[\tau_0,-T_g]$;
\item[\textup{(Q3)}] $|q(\tau)|\le\bar Q(1+\tau)^p$ for $\tau\in[0,\tau_1]$.
\end{description}
Put $\Errg=\int_{-T_g}^0e^{-x\tau}\err(\tau)\dd\tau$ and
$\Errpre=\int_{\tau_0}^{-T_g}e^{-x\tau}Q_{\mathrm{pre}}(\tau)\dd\tau$, and assume
\begin{description}[font=\normalfont]
\item[\textup{(Kick1)}] $\mu-1\le|\qs|/\big(10|z|\,\bar Q\,\Ip{p}{\sigma/4}\big)$;
\item[\textup{(Kick2)}] $e^{-\Re w\,T_g}\le\frac12$;
\item[\textup{(Kick3)}] $\Errg\le|\qs|/\big(20|z|(\mu-1)\big)$;
\item[\textup{(Kick4)}] $\Errpre\le|\qs|/\big(20|z|(\mu-1)\big)$.
\end{description}
Then at $\tau=\tau_1$,
\[
\frac{|\beta|\,|\qs|}{10|z|(\mu-1)}\le|A(\tau)|\le\frac{3|\beta|\,|\qs|\,\mu}{|z|(\mu-1)} .
\]
\end{proposition}

The hypotheses describe the drive $q$ on three consecutive intervals:
\textup{(Q2)} bounds it on $[\tau_0,-T_g]$, \textup{(Q1)} says that on $[-T_g,0]$
it is close to the exponential $\qs e^{z'\tau}$, and \textup{(Q3)} bounds its
growth on $[0,\tau_1]$. In Lemma~\ref{lem:kicktrue}, where $G'=G_{k+1}$ and
$T_g=\hat T_{k+1}$, these intervals are, for $k+1<n$, the quiet phase and the
active phase of ring $k+1$ and the pinning phase $\Pin_k$ of ring $k$, in kick
time; at the seeded level $k+1=n$ the proof of that lemma applies
Proposition~\ref{prop:kick} with a slightly smaller $T_g$. The proof below shows
that $\Re w>0$ and that $|w|$ is comparable to $|z|(\mu-1)/\mu$. By
\textup{(Q1)}, the integrand of Equation~\eqref{eq:kick-A} on $[-T_g,0]$ is
close to $\qs e^{ws}$. \textup{(Kick2)} makes $|1-e^{-wT_g}|\ge\frac12$, so
this interval contributes a term of modulus comparable to
$|\beta||\qs|/(|z|(\mu-1))$, and \textup{(Kick1)}, \textup{(Kick3)} and
\textup{(Kick4)} make the contributions of $[0,\tau_1]$, of the error in
\textup{(Q1)} and of $[\tau_0,-T_g]$ small against it.

\begin{proof}
We first compare $w$ with $z$ and bound $x$ from below. Put
$e_\mu=\frac{\mu-1}{8\mu}\le\frac1{16}$ and write $z'=z+d_1$,
$\lamst(G)=z+d_2$ with $|d_1|,|d_2|\le e_\mu\sigma$. Then
$w=z\frac{\mu-1}\mu+d_1-\frac{d_2}\mu$ with
$|d_1-\frac{d_2}\mu|\le2e_\mu\sigma=\frac{\sigma(\mu-1)}{4\mu}$, so, since
$\Re z=\sigma\le|z|$,
\[
\Re w\ge\frac{\sigma(\mu-1)}\mu-\frac{\sigma(\mu-1)}{4\mu}=\frac34\sigma\frac{\mu-1}\mu>0,\qquad
\frac34\,\frac{|z|(\mu-1)}\mu\le|w|\le\frac54\,\frac{|z|(\mu-1)}\mu .
\]
As $\mu\le2$, $x=\Re\lamst(G)/\mu\ge(1-e_\mu)\sigma/\mu\ge\frac{15}{32}\sigma\ge\frac\sigma4>0$,
and as $\Ip{p}{\cdot}$ is decreasing, $\Ip{p}{x}\le\Ip{p}{\sigma/4}$.

Next we split the integral of Equation~\eqref{eq:kick-A} at $-T_g$ and $0$,
according to the three hypotheses on $q$. Put $\mathsf K=\beta\qs(1-e^{-wT_g})/w$,
the main term, and let $\tau=\tau_1$. Then
\[
A(\tau)=\beta\int_{\tau_0}^{-T_g}e^{-\zeta_0s}q(s)\dd s+\beta\int_{-T_g}^0e^{-\zeta_0s}q(s)\dd s
+\beta\int_0^\tau e^{-\zeta_0s}q(s)\dd s ,
\]
and $|e^{-\zeta_0s}|=e^{-xs}$. By (Q2), the
first term on the right has modulus at most $|\beta|\Errpre$. Since
\[
\int_{-T_g}^0e^{-\zeta_0s}\qs e^{z's}\dd s=\qs\int_{-T_g}^0e^{ws}\dd s=\qs(1-e^{-wT_g})/w,
\]
the second term is $\mathsf K$ plus a remainder of modulus at most
$|\beta|\Errg$, by (Q1). By (Q3), $x>0$ and Equation~(I.\ref{I-eq:Ip}), the third
has modulus at most
$|\beta|\bar Q\int_0^\infty e^{-xs}(1+s)^p\dd s=|\beta|\bar Q\,\Ip{p}{x}$.

Finally we compare the main term and the three remainders with
$\mathsf Z=\frac25|\beta||\qs|/(|z|(\mu-1))$.
By (Kick2) and $|e^{-wT_g}|=e^{-\Re w\,T_g}$, $|1-e^{-wT_g}|\in[\frac12,\frac32]$,
and the bounds on $|w|$ from the first paragraph give
\[
|\mathsf K|\ge\frac12\cdot\frac45\cdot\frac{|\beta||\qs|\mu}{|z|(\mu-1)}=\mu\,\mathsf Z,\qquad
|\mathsf K|\le\frac32\cdot\frac43\cdot\frac{|\beta||\qs|\mu}{|z|(\mu-1)}=\frac{2|\beta||\qs|\mu}{|z|(\mu-1)} .
\]
By (Kick1) and $\Ip{p}{x}\le\Ip{p}{\sigma/4}$,
$|\beta|\bar Q\,\Ip{p}{x}\le|\beta||\qs|/(10|z|(\mu-1))=\frac14\mathsf Z$, and by
(Kick3) and (Kick4), $|\beta|\Errg,|\beta|\Errpre\le\frac18\mathsf Z$; hence
$|A(\tau)-\mathsf K|\le\frac12\mathsf Z$. Since $\mu>1$,
\begin{gather*}
|A(\tau)|\ge|\mathsf K|-\tfrac12\mathsf Z\ge\big(\mu-\tfrac12\big)\mathsf Z\ge\tfrac12\mathsf Z=\frac{|\beta||\qs|}{5|z|(\mu-1)},\\
|A(\tau)|\le|\mathsf K|+\tfrac12\mathsf Z\le\Big(2\mu+\frac15\Big)\frac{|\beta||\qs|}{|z|(\mu-1)}
\le\frac{3|\beta||\qs|\mu}{|z|(\mu-1)} .\qedhere
\end{gather*}
\end{proof}
